\section{Complete Experimental Specification}
\label{app:specification}
The primary experiment and follow-on experiment are separately recorded.
The primary grid uses $N=8$, depths $\{1,2,4,8\}$, candidate counts
$\{2,4,8\}$, active widths $\{1,2,4,8\}$, both snapshot endpoints, and
both graph families. Its seeds are 0--255 and each task is evaluated by
eight programmed controllers. The follow-on grids use the anchors
in the main-text ablation design; their full per-controller configurations
are in \path{results_extended/case_design.json}.

\paragraph{Pairing and randomness.}
The base generator draws a fixed eight-slot initial array, a permutation
of operation identifiers, a shuffled balanced query assignment, update
registers, nonzero increments, and candidate permutations. Candidate-count
levels filter the same eight-code permutation; active-width levels project
query and update indices into the active width. Across work levels, different
numbers of operations cause different subsequent random draws, so seed
pairing is not exact surface matching. Hidden resampling uses a separate
seeded random stream, draws eight fresh values after each success, and
replaces only active registers. Failed trials never advance that stream.
Controller sampling uses its own stream. Common controller seeds pair
variants, but differing controller logic can consume different numbers
of draws. These choices are intended for reproducibility, not to certify
statistical independence of every repeated condition.

\paragraph{Admission and positive controls.}
There is no outcome-based rejection or regeneration. All balanced tasks
must query every active register; clustered assignments preserve that
coverage and its multiplicities. Concentrated-query variants explicitly
have effective width one and are tagged as a degenerate manipulation.
Positive controls include unrestricted deterministic tracking, residue
tracking, deterministic preplanning, unrestricted tracking with a snapshot
after every hidden resampling, and candidate elimination with the worst-case
error allowance. The runner aborts if a designated positive control fails,
if the call horizon is exceeded, or if the adaptive cache exceeds two
entries. Analytic probabilities are derived only for supported cases;
unsupported predictions are labelled \texttt{not\_derived}, not replaced
with a default or fitted estimate.

\paragraph{Adaptive-cache details.}
Next-use-2 uses the public deterministic ready schedule to prioritize
registers queried in the future. On a visible snapshot, the current query
can be answered directly, so future-use selection skips that operation.
During history-only feedback, a successful action confirms the queried
residue from the chosen candidate code. The controller inserts that
residue, applies the public increment if relevant, and evicts values
whose next use is furthest away. This is not a generic learned memory
strategy; future-query metadata, successful-action feedback, and public
updates are deliberate advantages of the control.

\paragraph{Preparation and storage accounting.}
Residue caches store values modulo $b$; raw-value caches store values
modulo eight. Reported value bits exclude dictionary keys, graph metadata,
controller code, temporary operands, and runtime overhead. Preplanning
records $N$ candidate slots, using the initial state as preparation
workspace; indexed selection records $N$ maps of $b$ code-to-slot entries.
The preparation counter charges $Nb$ table entries for each of these
representations. It is a workload proxy, not a measurement of exact executed
CPU instructions or comparison count in Python's \texttt{index} function.
Lookup counts are kept separate from predicate comparisons. The experiment
does not aggregate these categories into an invented universal compute unit.

\section{Calibration Derivations}
\label{app:proofs}
\subsection{Completion states and schedule probabilities}
For $k=N/d$ disjoint chains, a completion set is determined by the completed
prefix length $0,\ldots,d$ of each chain, giving $(d+1)^k$ reachable masks.
Each chain imposes one internal order on its $d$ operations; interleaving
the $k$ orders gives $N!/(d!)^k$ valid permutations. A uniformly sampled
operation permutation therefore succeeds with probability $(d!)^{-k}$.

For barriers, all previous layers must be complete before any operation
in the next layer succeeds. Each of the $d$ layers contributes $2^k-1$
nonempty partial-completion choices, with one initial empty mask, giving
$d(2^k-1)+1$ distinct masks. Each layer's operations can be internally
permuted in $k!$ ways, so the valid-order count is $(k!)^d$ and random-order
success is $(k!)^d/N!$. Failed-action sinks are not included in these
completion-mask counts.

\subsection{Deterministic history and residue sufficiency}
Let $C_t$ be the completed-operation set after $t$ successes. Because each
update is additive and public,
\[
 x_r(t)=x_r(0)+
 \sum_{v\in C_t:u_v=r}\delta_v \pmod8.
\]
Thus the initial observation and successful-action history determine every
register exactly. Repeating snapshots cannot improve optimal information
availability in this deterministic protocol, although it can help a
restricted retention strategy. Since each tested $b$ divides eight, the
same equation modulo $b$ determines every candidate predicate. Tracking
$x_r\bmod b$ is therefore sufficient; storing all three raw value bits
is not necessary for $b=2$ or $b=4$.

A known ready schedule can also be simulated once using the initial
observation. Recording its future candidate slots is an exact plan in the
deterministic protocol. That construction establishes an alternative
representation, not a lower bound on memory and not a claim that planning
uses less total storage or preparation.

\subsection{Missing-information ceiling}
In the hidden-resampling condition, all active values are freshly uniform
after each success. Before a masked decision, conditioning on any visible
history still leaves the queried code uniform on $0,\ldots,b-1$.
Choosing a different ready operation cannot change this fact, because every
active register has the same independent fresh distribution. Consequently,
the conditional probability of each hidden decision succeeding is at most
$1/b$, regardless of additional computation.

With no retries, the number of visible decisions on a successful path is
$v=\lceil N/h\rceil$ for snapshot period $h>0$, or $v=1$ for initial-only
feedback. Applying the chain rule gives $P(\text{complete})\le b^{-(N-v)}$.
An unrestricted exact controller attains this ceiling. The bound does not
apply unchanged to recovery protocols: failed trials preserve state and
provide information that can support elimination. For scan-limited
controllers, formulas that average visible candidate-order draws are
ensemble predictions and are labelled separately from exact
instance-conditional predictions.

\subsection{Finite retry budgets}
Uniform replacement guesses have success probability $p=1/b$ per attempt.
If $F$ is the number of errors before $N$ successes, then
\[
 \Pr(F\le R)=\sum_{r=0}^{R}
   {N+r-1\choose r}p^N(1-p)^r.
\]
Under uniform elimination, the correct code has a uniformly random rank
in a permutation of the $b$ candidates. Errors per operation are therefore
independent uniform counts on $\{0,\ldots,b-1\}$. The finite-budget
probability is
\[
 \Pr(F\le R)=\frac{1}{b^N}
 \sum_{r=0}^{R}[z^r](1+z+\cdots+z^{b-1})^N.
\]
The implementation evaluates this coefficient sum by a truncated
convolution. If $R\ge N(b-1)$ the probability is one. Mean calls without
budget exhaustion are $Nb$ for replacement and $N(b+1)/2$ for elimination.
These are exact properties of the selected guessing policies, not
general guarantees for arbitrary tool-using controllers.

\section{All Follow-On Outcome Summaries}
\label{app:ablationtables}
Each row below averages equally over the remaining anchors of its grid.
Low and high are descriptive 95\% percentile bands from 2,000 paired
seed-block bootstrap resamples. They are not simultaneous confidence
statements. Constants produce degenerate bands; a sampled zero is not
a proof of zero probability. Calls include failed attempts and stop at
termination. For exact analytic predictions, progress, error counts,
comparison counts, preparation workload, lookups, and cache measurements,
see \path{results_extended/summary.csv}. All 77 paired contrast estimates
and their bands are in \path{results_extended/paired_contrasts.csv}.

Condition notation uses \texttt{n} for operation count, \texttt{p} for
snapshot period, \texttt{r} for global error allowance, and
\texttt{shock1} for independent hidden resampling. Period zero means
initial-only state feedback. The \texttt{reverse} flag indicates operation
identifier reversal. All grids use both graph families.

\small
\subsection{Operation Count}
\begin{longtable}{llrrrr}
\toprule
Condition & Controller & Success (\%) & Low & High & Calls\\
\midrule
\endhead
n16\_p0 & Cache-2 & 0.00 & 0.00 & 0.00 & 3.38\\
n16\_p0 & Exact & 100.00 & 100.00 & 100.00 & 16.00\\
n16\_p0 & Random schedule & 0.00 & 0.00 & 0.00 & 1.69\\
n16\_p0 & Scan-4 & 50.05 & 50.00 & 50.15 & 9.44\\
n16\_p1 & Cache-2 & 100.00 & 100.00 & 100.00 & 16.00\\
n16\_p1 & Exact & 100.00 & 100.00 & 100.00 & 16.00\\
n16\_p1 & Random schedule & 0.00 & 0.00 & 0.00 & 1.69\\
n16\_p1 & Scan-4 & 50.05 & 50.00 & 50.15 & 9.44\\
n4\_p0 & Cache-2 & 13.43 & 10.94 & 16.21 & 2.93\\
n4\_p0 & Exact & 100.00 & 100.00 & 100.00 & 4.00\\
n4\_p0 & Random schedule & 10.94 & 8.11 & 13.96 & 1.60\\
n4\_p0 & Scan-4 & 56.59 & 54.88 & 58.50 & 3.13\\
n4\_p1 & Cache-2 & 100.00 & 100.00 & 100.00 & 4.00\\
n4\_p1 & Exact & 100.00 & 100.00 & 100.00 & 4.00\\
n4\_p1 & Random schedule & 10.94 & 8.01 & 13.96 & 1.60\\
n4\_p1 & Scan-4 & 56.59 & 54.88 & 58.40 & 3.13\\
n8\_p0 & Cache-2 & 0.73 & 0.29 & 1.22 & 3.17\\
n8\_p0 & Exact & 100.00 & 100.00 & 100.00 & 8.00\\
n8\_p0 & Random schedule & 1.17 & 0.49 & 1.95 & 1.63\\
n8\_p0 & Scan-4 & 50.59 & 50.24 & 51.03 & 5.18\\
n8\_p1 & Cache-2 & 100.00 & 100.00 & 100.00 & 8.00\\
n8\_p1 & Exact & 100.00 & 100.00 & 100.00 & 8.00\\
n8\_p1 & Random schedule & 1.17 & 0.49 & 1.95 & 1.63\\
n8\_p1 & Scan-4 & 50.59 & 50.20 & 51.03 & 5.18\\
\bottomrule
\end{longtable}
\subsection{Retention Strategy}
\begin{longtable}{llrrrr}
\toprule
Condition & Controller & Success (\%) & Low & High & Calls\\
\midrule
\endhead
p0 & Next-use-2 & 0.05 & 0.00 & 0.15 & 5.05\\
p0 & Cache-2 & 0.00 & 0.00 & 0.00 & 2.98\\
p0 & Cache-4 & 50.00 & 50.00 & 50.00 & 9.66\\
p0 & Residue-8 & 100.00 & 100.00 & 100.00 & 16.00\\
p0 & Precomputed plan & 100.00 & 100.00 & 100.00 & 16.00\\
p0 & Stale-8 & 0.00 & 0.00 & 0.00 & 4.20\\
p4 & Next-use-2 & 9.52 & 7.86 & 11.28 & 7.68\\
p4 & Cache-2 & 0.00 & 0.00 & 0.00 & 3.10\\
p4 & Cache-4 & 50.00 & 50.00 & 50.00 & 9.76\\
p4 & Residue-8 & 100.00 & 100.00 & 100.00 & 16.00\\
p4 & Precomputed plan & 100.00 & 100.00 & 100.00 & 16.00\\
p4 & Stale-8 & 2.98 & 2.05 & 4.05 & 5.33\\
\bottomrule
\end{longtable}
\subsection{Snapshot Cadence and Hidden Resampling}
\begin{longtable}{llrrrr}
\toprule
Condition & Controller & Success (\%) & Low & High & Calls\\
\midrule
\endhead
p0\_shock0 & Cache-2 & 0.29 & 0.05 & 0.59 & 2.85\\
p0\_shock0 & Exact & 100.00 & 100.00 & 100.00 & 8.00\\
p0\_shock0 & Scan-4 & 50.68 & 50.34 & 51.12 & 5.18\\
p0\_shock0 & Uniform & 0.00 & 0.00 & 0.00 & 1.25\\
p0\_shock1 & Cache-2 & 0.00 & 0.00 & 0.00 & 2.25\\
p0\_shock1 & Exact & 0.00 & 0.00 & 0.00 & 2.25\\
p0\_shock1 & Scan-4 & 0.00 & 0.00 & 0.00 & 2.01\\
p0\_shock1 & Uniform & 0.00 & 0.00 & 0.00 & 1.25\\
p1\_shock0 & Cache-2 & 100.00 & 100.00 & 100.00 & 8.00\\
p1\_shock0 & Exact & 100.00 & 100.00 & 100.00 & 8.00\\
p1\_shock0 & Scan-4 & 50.68 & 50.34 & 51.12 & 5.18\\
p1\_shock0 & Uniform & 0.00 & 0.00 & 0.00 & 1.25\\
p1\_shock1 & Cache-2 & 100.00 & 100.00 & 100.00 & 8.00\\
p1\_shock1 & Exact & 100.00 & 100.00 & 100.00 & 8.00\\
p1\_shock1 & Scan-4 & 50.98 & 50.54 & 51.46 & 5.25\\
p1\_shock1 & Uniform & 0.00 & 0.00 & 0.00 & 1.25\\
p2\_shock0 & Cache-2 & 4.64 & 3.66 & 5.71 & 3.58\\
p2\_shock0 & Exact & 100.00 & 100.00 & 100.00 & 8.00\\
p2\_shock0 & Scan-4 & 50.68 & 50.34 & 51.17 & 5.18\\
p2\_shock0 & Uniform & 0.00 & 0.00 & 0.00 & 1.25\\
p2\_shock1 & Cache-2 & 0.29 & 0.05 & 0.59 & 2.49\\
p2\_shock1 & Exact & 0.29 & 0.05 & 0.59 & 2.49\\
p2\_shock1 & Scan-4 & 0.24 & 0.05 & 0.54 & 2.21\\
p2\_shock1 & Uniform & 0.00 & 0.00 & 0.00 & 1.25\\
p4\_shock0 & Cache-2 & 0.44 & 0.20 & 0.73 & 2.93\\
p4\_shock0 & Exact & 100.00 & 100.00 & 100.00 & 8.00\\
p4\_shock0 & Scan-4 & 50.68 & 50.29 & 51.07 & 5.18\\
p4\_shock0 & Uniform & 0.00 & 0.00 & 0.00 & 1.25\\
p4\_shock1 & Cache-2 & 0.00 & 0.00 & 0.00 & 2.27\\
p4\_shock1 & Exact & 0.00 & 0.00 & 0.00 & 2.27\\
p4\_shock1 & Scan-4 & 0.00 & 0.00 & 0.00 & 2.03\\
p4\_shock1 & Uniform & 0.00 & 0.00 & 0.00 & 1.25\\
\bottomrule
\end{longtable}
\subsection{Recovery and Candidate Elimination}
\begin{longtable}{llrrrr}
\toprule
Condition & Controller & Success (\%) & Low & High & Calls\\
\midrule
\endhead
r0\_p0 & Cache-2 + eliminate & 0.29 & 0.05 & 0.59 & 2.85\\
r0\_p0 & Scan-1 + eliminate & 0.29 & 0.10 & 0.54 & 1.64\\
r0\_p0 & Uniform & 0.00 & 0.00 & 0.00 & 1.25\\
r0\_p0 & Eliminate & 0.00 & 0.00 & 0.00 & 1.25\\
r0\_p1 & Cache-2 + eliminate & 100.00 & 100.00 & 100.00 & 8.00\\
r0\_p1 & Scan-1 + eliminate & 0.29 & 0.10 & 0.54 & 1.64\\
r0\_p1 & Uniform & 0.00 & 0.00 & 0.00 & 1.25\\
r0\_p1 & Eliminate & 0.00 & 0.00 & 0.00 & 1.25\\
r2\_p0 & Cache-2 + eliminate & 3.52 & 2.59 & 4.49 & 5.86\\
r2\_p0 & Scan-1 + eliminate & 2.93 & 2.10 & 3.76 & 5.11\\
r2\_p0 & Uniform & 0.05 & 0.00 & 0.15 & 3.72\\
r2\_p0 & Eliminate & 0.00 & 0.00 & 0.00 & 3.92\\
r2\_p1 & Cache-2 + eliminate & 100.00 & 100.00 & 100.00 & 8.00\\
r2\_p1 & Scan-1 + eliminate & 2.93 & 2.10 & 3.81 & 5.11\\
r2\_p1 & Uniform & 0.05 & 0.00 & 0.15 & 3.72\\
r2\_p1 & Eliminate & 0.00 & 0.00 & 0.00 & 3.92\\
r56\_p0 & Cache-2 + eliminate & 100.00 & 100.00 & 100.00 & 18.94\\
r56\_p0 & Scan-1 + eliminate & 100.00 & 100.00 & 100.00 & 21.52\\
r56\_p0 & Uniform & 78.56 & 76.61 & 80.52 & 43.27\\
r56\_p0 & Eliminate & 100.00 & 100.00 & 100.00 & 27.87\\
r56\_p1 & Cache-2 + eliminate & 100.00 & 100.00 & 100.00 & 8.00\\
r56\_p1 & Scan-1 + eliminate & 100.00 & 100.00 & 100.00 & 21.52\\
r56\_p1 & Uniform & 78.56 & 76.71 & 80.47 & 43.27\\
r56\_p1 & Eliminate & 100.00 & 100.00 & 100.00 & 27.87\\
r8\_p0 & Cache-2 + eliminate & 44.63 & 42.77 & 46.39 & 13.63\\
r8\_p0 & Scan-1 + eliminate & 43.07 & 41.75 & 44.38 & 12.81\\
r8\_p0 & Uniform & 1.27 & 0.78 & 1.81 & 11.09\\
r8\_p0 & Eliminate & 6.54 & 5.42 & 7.67 & 12.64\\
r8\_p1 & Cache-2 + eliminate & 100.00 & 100.00 & 100.00 & 8.00\\
r8\_p1 & Scan-1 + eliminate & 43.07 & 41.75 & 44.43 & 12.81\\
r8\_p1 & Uniform & 1.27 & 0.78 & 1.81 & 11.09\\
r8\_p1 & Eliminate & 6.54 & 5.42 & 7.67 & 12.64\\
\bottomrule
\end{longtable}
\subsection{Query Locality and State Updates}
\begin{longtable}{llrrrr}
\toprule
Condition & Controller & Success (\%) & Low & High & Calls\\
\midrule
\endhead
balanced\_additive & Next-use-2 & 0.05 & 0.00 & 0.15 & 5.05\\
balanced\_additive & Cache-2 & 0.00 & 0.00 & 0.00 & 2.98\\
balanced\_additive & Exact & 100.00 & 100.00 & 100.00 & 16.00\\
balanced\_additive & Stale-8 & 0.00 & 0.00 & 0.00 & 4.20\\
balanced\_static & Next-use-2 & 0.05 & 0.00 & 0.15 & 5.05\\
balanced\_static & Cache-2 & 0.00 & 0.00 & 0.00 & 3.04\\
balanced\_static & Exact & 100.00 & 100.00 & 100.00 & 16.00\\
balanced\_static & Stale-8 & 100.00 & 100.00 & 100.00 & 16.00\\
clustered\_additive & Next-use-2 & 1.86 & 1.32 & 2.44 & 7.67\\
clustered\_additive & Cache-2 & 0.00 & 0.00 & 0.00 & 7.26\\
clustered\_additive & Exact & 100.00 & 100.00 & 100.00 & 16.00\\
clustered\_additive & Stale-8 & 0.15 & 0.00 & 0.39 & 4.80\\
clustered\_static & Next-use-2 & 1.76 & 1.17 & 2.34 & 7.68\\
clustered\_static & Cache-2 & 0.00 & 0.00 & 0.00 & 7.24\\
clustered\_static & Exact & 100.00 & 100.00 & 100.00 & 16.00\\
clustered\_static & Stale-8 & 100.00 & 100.00 & 100.00 & 16.00\\
concentrated\_additive & Next-use-2 & 100.00 & 100.00 & 100.00 & 16.00\\
concentrated\_additive & Cache-2 & 100.00 & 100.00 & 100.00 & 16.00\\
concentrated\_additive & Exact & 100.00 & 100.00 & 100.00 & 16.00\\
concentrated\_additive & Stale-8 & 9.42 & 6.93 & 12.11 & 7.21\\
concentrated\_static & Next-use-2 & 100.00 & 100.00 & 100.00 & 16.00\\
concentrated\_static & Cache-2 & 100.00 & 100.00 & 100.00 & 16.00\\
concentrated\_static & Exact & 100.00 & 100.00 & 100.00 & 16.00\\
concentrated\_static & Stale-8 & 100.00 & 100.00 & 100.00 & 16.00\\
\bottomrule
\end{longtable}
\subsection{Candidate Selection and Resource Controls}
\begin{longtable}{llrrrr}
\toprule
Condition & Controller & Success (\%) & Low & High & Calls\\
\midrule
\endhead
b2 & Exact & 100.00 & 100.00 & 100.00 & 8.00\\
b2 & Indexed & 100.00 & 100.00 & 100.00 & 8.00\\
b2 & Scan-1 & 100.00 & 100.00 & 100.00 & 8.00\\
b2 & Sham-check-4 & 100.00 & 100.00 & 100.00 & 8.00\\
b2 & Scan-2 & 100.00 & 100.00 & 100.00 & 8.00\\
b2 & Fixed-check-4 & 100.00 & 100.00 & 100.00 & 8.00\\
b4 & Exact & 100.00 & 100.00 & 100.00 & 8.00\\
b4 & Indexed & 100.00 & 100.00 & 100.00 & 8.00\\
b4 & Scan-1 & 0.59 & 0.00 & 1.37 & 2.00\\
b4 & Sham-check-4 & 0.59 & 0.00 & 1.37 & 2.00\\
b4 & Scan-2 & 9.18 & 6.64 & 12.11 & 3.55\\
b4 & Fixed-check-4 & 100.00 & 100.00 & 100.00 & 8.00\\
b8 & Exact & 100.00 & 100.00 & 100.00 & 8.00\\
b8 & Indexed & 100.00 & 100.00 & 100.00 & 8.00\\
b8 & Scan-1 & 0.00 & 0.00 & 0.00 & 1.25\\
b8 & Sham-check-4 & 0.00 & 0.00 & 0.00 & 1.25\\
b8 & Scan-2 & 0.00 & 0.00 & 0.00 & 1.61\\
b8 & Fixed-check-4 & 1.17 & 0.20 & 2.34 & 2.39\\
\bottomrule
\end{longtable}
\subsection{Operation Identifiers}
\begin{longtable}{llrrrr}
\toprule
Condition & Controller & Success (\%) & Low & High & Calls\\
\midrule
\endhead
reverse0\_p0 & Cache-2 & 0.29 & 0.05 & 0.59 & 2.85\\
reverse0\_p0 & Exact & 100.00 & 100.00 & 100.00 & 8.00\\
reverse0\_p0 & Random schedule & 0.00 & 0.00 & 0.00 & 1.29\\
reverse0\_p0 & Scan-4 & 50.68 & 50.34 & 51.12 & 5.18\\
reverse0\_p1 & Cache-2 & 100.00 & 100.00 & 100.00 & 8.00\\
reverse0\_p1 & Exact & 100.00 & 100.00 & 100.00 & 8.00\\
reverse0\_p1 & Random schedule & 0.00 & 0.00 & 0.00 & 1.29\\
reverse0\_p1 & Scan-4 & 50.68 & 50.29 & 51.12 & 5.18\\
reverse1\_p0 & Cache-2 & 0.15 & 0.00 & 0.34 & 2.88\\
reverse1\_p0 & Exact & 100.00 & 100.00 & 100.00 & 8.00\\
reverse1\_p0 & Random schedule & 0.00 & 0.00 & 0.00 & 1.35\\
reverse1\_p0 & Scan-4 & 50.44 & 50.15 & 50.83 & 5.23\\
reverse1\_p1 & Cache-2 & 100.00 & 100.00 & 100.00 & 8.00\\
reverse1\_p1 & Exact & 100.00 & 100.00 & 100.00 & 8.00\\
reverse1\_p1 & Random schedule & 0.00 & 0.00 & 0.00 & 1.35\\
reverse1\_p1 & Scan-4 & 50.44 & 50.10 & 50.83 & 5.23\\
\bottomrule
\end{longtable}
\normalsize

\section{All Endpoint Interaction Estimates}
\label{app:interactions}
The following estimates are exploratory follow-on analyses of the original
192-cell factorial experiment. $D,B,W,O$ denote depth, candidate count,
active width, and history-only versus snapshot feedback. Endpoint
contrasts use high minus low for each participating dimension and average
all nonparticipating controls over their original levels. A pairwise
interaction is a difference-in-differences; the corresponding signed
construction generalizes to orders three and four. Values are percentage
points, so an interaction is not itself a success probability. Per-estimate
bootstrap bands are retained in
\path{results_extended/all_endpoint_interactions.csv}.

\begin{center}
\small
\resizebox{\textwidth}{!}{\begin{tabular}{llrrrrrr}
\toprule
Family & Controller & $D\times{} B$ & $D\times{} W$ & $D\times{} O$ & $B\times{} W$ & $B\times{} O$ & $W\times{} O$\\
\midrule
Barriers & exact & 0.00 & 0.00 & 0.00 & 0.00 & 0.00 & 0.00\\
Barriers & memory2 & -0.10 & 0.10 & 0.10 & -2.05 & -2.05 & -97.95\\
Barriers & memory2\_scan4 & 0.49 & -1.27 & 0.10 & 48.93 & 47.56 & -50.68\\
Barriers & memory4 & 0.68 & -0.68 & -0.68 & -4.00 & -4.00 & -95.61\\
Barriers & random\_order & 0.00 & 0.00 & 0.00 & 0.00 & 0.00 & 0.00\\
Barriers & scan1 & 0.00 & 0.00 & 0.00 & 0.00 & 0.00 & 0.00\\
Barriers & scan4 & -0.59 & 0.59 & 0.00 & 2.93 & 0.00 & 0.00\\
Barriers & uniform & -0.59 & 0.98 & 0.00 & -0.59 & 0.00 & 0.00\\
Chains & exact & 0.00 & 0.00 & 0.00 & 0.00 & 0.00 & 0.00\\
Chains & memory2 & -0.10 & 0.10 & 0.10 & -2.05 & -2.05 & -97.95\\
Chains & memory2\_scan4 & 0.49 & -1.27 & 0.10 & 48.93 & 47.56 & -50.68\\
Chains & memory4 & 0.68 & -0.68 & -0.68 & -4.00 & -4.00 & -95.61\\
Chains & random\_order & 0.00 & 0.00 & 0.00 & 0.00 & 0.00 & 0.00\\
Chains & scan1 & 0.00 & 0.00 & 0.00 & 0.00 & 0.00 & 0.00\\
Chains & scan4 & -0.59 & 0.59 & 0.00 & 2.93 & 0.00 & 0.00\\
Chains & uniform & -0.59 & 0.98 & 0.00 & -0.59 & 0.00 & 0.00\\
\bottomrule
\end{tabular}}

\medskip
\resizebox{\textwidth}{!}{\begin{tabular}{llrrrrr}
\toprule
Family & Controller & $D\times{} B\times{} W$ & $D\times{} B\times{} O$ & $D\times{} W\times{} O$ & $B\times{} W\times{} O$ & $D\times{} B\times{} W\times{} O$\\
\midrule
Barriers & exact & 0.00 & 0.00 & 0.00 & 0.00 & 0.00\\
Barriers & memory2 & -0.20 & -0.20 & 0.20 & -4.10 & -0.39\\
Barriers & memory2\_scan4 & -1.37 & 1.37 & 0.20 & 95.12 & 2.73\\
Barriers & memory4 & 1.37 & 1.37 & -1.37 & -8.01 & 2.73\\
Barriers & random\_order & 0.00 & 0.00 & 0.00 & 0.00 & 0.00\\
Barriers & scan1 & 0.00 & 0.00 & 0.00 & 0.00 & 0.00\\
Barriers & scan4 & 1.17 & 0.00 & 0.00 & 0.00 & 0.00\\
Barriers & uniform & -1.95 & 0.00 & 0.00 & 0.00 & 0.00\\
Chains & exact & 0.00 & 0.00 & 0.00 & 0.00 & 0.00\\
Chains & memory2 & -0.20 & -0.20 & 0.20 & -4.10 & -0.39\\
Chains & memory2\_scan4 & -1.37 & 1.37 & 0.20 & 95.12 & 2.73\\
Chains & memory4 & 1.37 & 1.37 & -1.37 & -8.01 & 2.73\\
Chains & random\_order & 0.00 & 0.00 & 0.00 & 0.00 & 0.00\\
Chains & scan1 & 0.00 & 0.00 & 0.00 & 0.00 & 0.00\\
Chains & scan4 & 1.17 & 0.00 & 0.00 & 0.00 & 0.00\\
Chains & uniform & -1.95 & 0.00 & 0.00 & 0.00 & 0.00\\
\bottomrule
\end{tabular}}
\end{center}

\section{Workshop Comparison and Submission Scope}
We inspected four ALA 2025 papers as a purposive sample, not all workshop
papers and not a systematic novelty search. Their total PDF lengths are
nine pages (failure sampling), ten (Dynamic Thinker), ten (task-graph
scheduling), and eight (deliberation cost). References begin on pages
nine, eight, seven, and eight, respectively. Total PDF length therefore
does not equal the allowed main-text length.

The historical ALA 2025 call allows original work up to eight pages
excluding references, and preliminary or outlook papers up to six pages.
It specifies double-blind review and deanonymized camera-ready copies.
Our expanded authored version is a preprint-style manuscript with an
appendix; the separate \path{review.tex} source removes the author block
and appendix, but must still be checked against the current workshop's
track and page rules. A longer document is not evidence of acceptance
or of a stronger contribution. No current-edition acceptance, conference
metadata, or submission date is asserted.

\section{Reproduction and Evidence Boundaries}
The benchmark, controllers, tests, and analysis use the Python standard
library. Document inspection additionally uses the pinned PDF packages.
From the project directory in a Windows PowerShell terminal:
\begin{verbatim}
.\.venv\Scripts\python.exe -m unittest -v test_benchmark test_extended
.\.venv\Scripts\python.exe run_experiments.py --seeds 256
.\.venv\Scripts\python.exe analyze.py
.\.venv\Scripts\python.exe run_extended.py
.\.venv\Scripts\python.exe analyze_extended.py
.\build.ps1
.\.venv\Scripts\python.exe verify_artifact.py
.\.venv\Scripts\python.exe package_artifact.py
\end{verbatim}
Re-running the primary experiment replaces its generated records;
the source and protocol hashes must match the supplied metadata when
reproducing that study. The original and follow-on episode files have
separate hashes and execution metadata. The two locally frozen protocols
and the time ordering of exploratory analyses are retained explicitly.
Complete primary contrasts, rare-scheduler Wilson intervals, and example
traces remain available in the original result directory.

This artifact contains only programmed-controller evidence. No frontier
model, learned policy, training run, provider reasoning-budget intervention,
human participant, or realistic transfer task has been evaluated.
Source compilation and local reproducibility do not establish arXiv
moderation approval, workshop acceptance, or compatibility with a
particular submission processor. The author must review the manuscript
and current submission rules before uploading it.
